\begin{proof}[Proof of \cref{obj:prop:bounded-submodel-comparison}]
Fix the admissible tuple \(v=(p,\alpha,\beta,\gamma)\) and its matched smoothness tuple \(w=(\alpha,\beta,\gamma)\).

\begin{enumerate}
\item We first establish the integrated conversion identity. Write
\[
\mathcal O=[0,1]\times\{0,1\}\times\mathbb R
\]
for the Borel record space. For a law \(P\) in either
\(\mathcal M_v\) or \(\mathcal M_{(2,\alpha,\beta,\gamma)}\), let
\(Q_{a_{\mathrm{arm}},P}(\,\cdot\mid x)\) be its conditional outcome probability kernel in arm \(a_{\mathrm{arm}}\in\{0,1\}\). Thus
\[
\begin{split}
P(dx,\{a_{\mathrm{arm}}\},dy)
&=\lambda(dx)e_P(x)^{a_{\mathrm{arm}}}
  (1-e_P(x))^{1-a_{\mathrm{arm}}}
  Q_{a_{\mathrm{arm}},P}(dy\mid x),\\
m_{a_{\mathrm{arm}},P}(x)
&=m_{0,P}(x)+a_{\mathrm{arm}}\tau_P(x)
 =\int y\,Q_{a_{\mathrm{arm}},P}(dy\mid x)
 \qquad\text{for \(\lambda\)-almost every \(x\)}.
\end{split}
\]
The mean bounds give
\[
|m_{0,P}(x)|\le\frac12,\qquad
|m_{1,P}(x)|
\le |m_{0,P}(x)|+|\tau_P(x)|\le1.
\]
Define the signed-binary conversion \(\mathcal C_{\mathrm{bin}}(P)\) by
\[
\begin{split}
Q^{\mathrm{bin}}_{a_{\mathrm{arm}},P}(dy\mid x)
&=\frac{1+m_{a_{\mathrm{arm}},P}(x)}2\,\delta_1(dy)
 +\frac{1-m_{a_{\mathrm{arm}},P}(x)}2\,\delta_{-1}(dy),\\
\mathcal C_{\mathrm{bin}}(P)(dx,\{a_{\mathrm{arm}}\},dy)
&=\lambda(dx)e_P(x)^{a_{\mathrm{arm}}}
 (1-e_P(x))^{1-a_{\mathrm{arm}}}
 Q^{\mathrm{bin}}_{a_{\mathrm{arm}},P}(dy\mid x).
\end{split}
\]
The two arm weights are nonnegative and sum to one, and
\[
\int y\,Q^{\mathrm{bin}}_{a_{\mathrm{arm}},P}(dy\mid x)
=m_{a_{\mathrm{arm}},P}(x),\qquad
\int |y|^2\,Q^{\mathrm{bin}}_{a_{\mathrm{arm}},P}(dy\mid x)=1.
\]
Consequently the conversion has signed-binary support and retains the uniform design, overlap, smoothness restrictions, and mean bounds. In particular,
\[
\begin{gathered}
\mathcal C_{\mathrm{bin}}(P)\in\mathcal M_w^{\mathrm{bin}},\qquad
\bigl(\mathcal C_{\mathrm{bin}}(P)\bigr)_{X,A}=P_{X,A},\\
e_{\mathcal C_{\mathrm{bin}}(P)}=e_P,\qquad
m_{0,\mathcal C_{\mathrm{bin}}(P)}=m_{0,P},\qquad
\tau_{\mathcal C_{\mathrm{bin}}(P)}=\tau_P.
\end{gathered}
\]
Substitution into the definitions of distance and constancy yields
\[
d(\mathcal C_{\mathrm{bin}}(P))=d(P),
\]
and
\[
\begin{split}
&\exists c\in[-1/2,1/2]\ 
  \forall x,\ \tau_{\mathcal C_{\mathrm{bin}}(P)}(x)=c\\
&\hspace{35mm}\Longleftrightarrow
\exists c\in[-1/2,1/2]\ 
  \forall x,\ \tau_P(x)=c.
\end{split}
\]
% lean: CausalSmith.Stat.FinitepHomogeneityDensegamma.binaryConversion_primitives
% lean: CausalSmith.Stat.FinitepHomogeneityDensegamma.binaryConversion_inModel
% lean: CausalSmith.Stat.FinitepHomogeneityDensegamma.binaryConversion_preserves_effect

To define record randomization on the whole record space, set
\[
\ell_1(y)=\max\{-1,\min\{1,y\}\}\qquad(y\in\mathbb R)
\]
and define the probability kernel
\[
\mathsf K_{\mathrm{sign}}((x,a_{\mathrm{obs}},y),\cdot)
=
\frac{1+\ell_1(y)}2\,\delta_{(x,a_{\mathrm{obs}},1)}
+\frac{1-\ell_1(y)}2\,\delta_{(x,a_{\mathrm{obs}},-1)},
\qquad (x,a_{\mathrm{obs}},y)\in\mathcal O.
\]
For \(P\in\mathcal M_w^{\mathrm b}\), uniform design and fixed overlap transfer the observed outcome bound to both conditional arm kernels:
\[
Q_{a_{\mathrm{arm}},P}([-1,1]\mid x)=1
\qquad
\text{for \(\lambda\)-almost every \(x\), for both arms}.
\]
Indeed, disintegrating \(P(|Y|>1)=0\) gives
\[
\int
\left[
(1-e_P(x))Q_{0,P}(\{|y|>1\}\mid x)
+e_P(x)Q_{1,P}(\{|y|>1\}\mid x)
\right]\lambda(dx)=0.
\]
Both summands are nonnegative, and both arm weights are at least \(1/4\). Each conditional probability therefore vanishes almost everywhere. Since \(\ell_1(y)=y\) on \([-1,1]\), for \(b_{\mathrm{sign}}\in\{-1,1\}\),
\[
\int
\frac{1+b_{\mathrm{sign}}\ell_1(y)}2\,
Q_{a_{\mathrm{arm}},P}(dy\mid x)
=
\frac{1+b_{\mathrm{sign}}m_{a_{\mathrm{arm}},P}(x)}2
\qquad\text{for \(\lambda\)-almost every \(x\)}.
\]
Integrating the two atomic weights and then the treatment mixture proves the one-record identity
\[
\int_{\mathcal O}\mathsf K_{\mathrm{sign}}(o,\cdot)\,P(do)
=\mathcal C_{\mathrm{bin}}(P).
\]
% lean: CausalSmith.Stat.FinitepHomogeneityDensegamma.binaryRecordKernel_comp

Now let \(n\in\mathbb N\), let \(\phi\) be a measurable \([0,1]\)-valued test, and write its input as
\[
\boldsymbol o=(o_i)_{i=1}^n\in\mathcal O^n,\qquad
o_i=(x_i,a_{\mathrm{obs},i},y_i),\qquad
u_{\mathrm{seed}}\in[0,1].
\]
Define its integrated rejection map by
\[
\widetilde\phi(\boldsymbol o,u_{\mathrm{seed}})
=
\sum_{\boldsymbol b_{\mathrm{sign}}\in\{-1,1\}^n}
\phi\!\left(
 ((x_i,a_{\mathrm{obs},i},b_{\mathrm{sign},i}))_{i=1}^n,
 u_{\mathrm{seed}}
\right)
\prod_{i=1}^n
\frac{1+b_{\mathrm{sign},i}\ell_1(y_i)}2,
\qquad
\boldsymbol b_{\mathrm{sign}}
=(b_{\mathrm{sign},i})_{i=1}^n.
\]
Its weights are nonnegative and satisfy
\[
\sum_{\boldsymbol b_{\mathrm{sign}}\in\{-1,1\}^n}
\prod_{i=1}^n
\frac{1+b_{\mathrm{sign},i}\ell_1(y_i)}2
=
\prod_{i=1}^n
\left(
\frac{1+\ell_1(y_i)}2+\frac{1-\ell_1(y_i)}2
\right)=1.
\]
Thus \(\widetilde\phi\) is measurable and takes values in \([0,1]\). For \(n=0\), the sign-vector set contains the single empty vector and the empty product is one, so the same definition applies.

The finite sum is the integral of \(\phi\) against the independent product of the record kernels. Independence and the one-record identity imply
\[
\int_{\mathcal O^n}
\bigotimes_{i=1}^n\mathsf K_{\mathrm{sign}}(o_i,\cdot)\,
P^{\otimes n}(d\boldsymbol o)
=
\mathcal C_{\mathrm{bin}}(P)^{\otimes n}.
\]
This follows by factorization on measurable rectangles; the empty product gives the identity for \(n=0\). Hence, for every fixed \(u_{\mathrm{seed}}\),
\[
\int_{\mathcal O^n}
\widetilde\phi(\boldsymbol o,u_{\mathrm{seed}})
\,P^{\otimes n}(d\boldsymbol o)
=
\int_{\mathcal O^n}
\phi(\boldsymbol o,u_{\mathrm{seed}})
\,\mathcal C_{\mathrm{bin}}(P)^{\otimes n}(d\boldsymbol o).
\]
Integrating in the public seed and using the experiment of
\cref{obj:def:original-record-experiment} gives
\[
\mathsf R_n(P,\widetilde\phi)
=
\mathsf R_n(\mathcal C_{\mathrm{bin}}(P),\phi).
\]
Boundedness of the tests justifies the integral exchanges. On bounded records, the kernel uses exactly the probabilities \((1+Y)/2\) and \((1-Y)/2\) stated in the proposition.
% lean: CausalSmith.Stat.FinitepHomogeneityDensegamma.binaryRandomization_rejectProb

\item We next compare the binary and bounded testing risks. In the infimum--supremum comparisons, a supremum of type-II errors over an empty alternative set is assigned the value zero. The test
\[
\phi_{\mathrm{zero}}(\boldsymbol o,u_{\mathrm{seed}})=0
\]
is level-valid for every null class. The sets of level-valid tests are therefore nonempty, and
\[
0\le 1-\mathsf R_n(P,\phi)\le1
\]
bounds every individual type-II error and every corresponding worst-case error.

Signed-binary support implies bounded support. Keeping the remaining restrictions and the constancy condition gives
\[
\mathcal M_w^{\mathrm{bin}}\subseteq\mathcal M_w^{\mathrm b},
\qquad
H_0^{\mathrm{bin}}(w)\subseteq H_0^{\mathrm b}(w).
\]
Every bounded-null level-valid test is consequently binary-null level-valid, with
\[
\sup_{\substack{P\in\mathcal M_w^{\mathrm{bin}}\\d(P)\ge r}}
\{1-\mathsf R_n(P,\phi)\}
\le
\sup_{\substack{P\in\mathcal M_w^{\mathrm b}\\d(P)\ge r}}
\{1-\mathsf R_n(P,\phi)\}.
\]
For a nonempty binary alternative set, this is the supremum comparison under inclusion; for an empty binary alternative set, it follows from the zero convention and nonnegativity. Taking infima gives
\[
B_n^{\mathrm{bin}}(w,r)\le B_n^{\mathrm b}(w,r).
\]
% lean: CausalSmith.Stat.FinitepHomogeneityDensegamma.binaryTestingRisk_le_boundedTestingRisk

Conversely, take any binary-null level-valid test \(\phi\). For every \(P\in H_0^{\mathrm b}(w)\), the conversion belongs to \(H_0^{\mathrm{bin}}(w)\), so the preceding identity gives
\[
\mathsf R_n(P,\widetilde\phi)
=
\mathsf R_n(\mathcal C_{\mathrm{bin}}(P),\phi)
\le\frac1{10}.
\]
For a bounded alternative \(P\) with \(d(P)\ge r\), its conversion is a binary alternative with the same distance, and
\[
1-\mathsf R_n(P,\widetilde\phi)
=
1-\mathsf R_n(\mathcal C_{\mathrm{bin}}(P),\phi)
\le
\sup_{\substack{P\in\mathcal M_w^{\mathrm{bin}}\\d(P)\ge r}}
\{1-\mathsf R_n(P,\phi)\}.
\]
If the bounded alternative set is nonempty, taking its supremum preserves this inequality; the binary supremum is bounded above by one. If the bounded alternative set is empty, its worst-case error is zero and the binary worst-case error is nonnegative. Thus in either case,
\[
\sup_{\substack{P\in\mathcal M_w^{\mathrm b}\\d(P)\ge r}}
\{1-\mathsf R_n(P,\widetilde\phi)\}
\le
\sup_{\substack{P\in\mathcal M_w^{\mathrm{bin}}\\d(P)\ge r}}
\{1-\mathsf R_n(P,\phi)\}.
\]
The bounded-model infimum is at most the left-hand side for each such \(\phi\). Taking the binary-model infimum proves
\[
B_n^{\mathrm b}(w,r)\le B_n^{\mathrm{bin}}(w,r),
\qquad
B_n^{\mathrm{bin}}(w,r)=B_n^{\mathrm b}(w,r).
\]
% lean: CausalSmith.Stat.FinitepHomogeneityDensegamma.boundedTestingRisk_eq_binary_of_conversion

\item For the remaining model inclusion, take \(P\in\mathcal M_w^{\mathrm b}\). The conditional support established above and the nonnegative moment exponent supplied by admissibility give
\[
\int |y|^p\,Q_{a_{\mathrm{arm}},P}(dy\mid x)
\le
\int 1\,Q_{a_{\mathrm{arm}},P}(dy\mid x)
=1\le10
\qquad\text{for \(\lambda\)-almost every \(x\)}.
\]
All the other primitive restrictions are identical at the matched smoothness tuple. Therefore
\[
\mathcal M_w^{\mathrm{bin}}
\subseteq\mathcal M_w^{\mathrm b}
\subseteq\mathcal M_v,
\qquad
H_0^{\mathrm b}(w)\subseteq H_0(v).
\]
% lean: CausalSmith.Stat.FinitepHomogeneityDensegamma.boundedModel_inModel

\item The distance suprema are equal because the attained distance sets are equal. The binary-to-bounded inclusion gives one direction, and converting each bounded law gives the other:
\[
\{d(P):P\in\mathcal M_w^{\mathrm{bin}}\}
=
\{d(P):P\in\mathcal M_w^{\mathrm b}\}.
\]
For comparison with the full model, convert each bounded law. Its converted arm kernels also satisfy
\[
\int |y|^p\,Q^{\mathrm{bin}}_{a_{\mathrm{arm}},P}(dy\mid x)
=1\le10,
\]
so the conversion belongs to \(\mathcal M_v\) and has the same distance. Conversely, the conversion of every \(P\in\mathcal M_v\) belongs to the binary model, hence to the bounded model, and retains \(d(P)\). Thus
\[
\{d(P):P\in\mathcal M_w^{\mathrm b}\}
=
\{d(P):P\in\mathcal M_v\}.
\]
Taking suprema yields
\[
D_w^{\mathrm{bin}}=D_w^{\mathrm b}=D_v.
\]
Finally, \cref{obj:lem:causal-nonempty} supplies the boundedness of the full-model distance set and the witness lower bound
\[
d_0\le D_v\le\frac12.
\]
% lean: CausalSmith.Stat.FinitepHomogeneityDensegamma.maxDistBinary_eq_maxDistBounded
% lean: CausalSmith.Stat.FinitepHomogeneityDensegamma.maxDistBounded_eq_maxDist
% lean: CausalSmith.Stat.FinitepHomogeneityDensegamma.model_distance_lower

\item A test level-valid for \(H_0(v)\) is level-valid for \(H_0^{\mathrm b}(w)\), and model inclusion gives
\[
\sup_{\substack{P\in\mathcal M_w^{\mathrm b}\\d(P)\ge r}}
\{1-\mathsf R_n(P,\phi)\}
\le
\sup_{\substack{P\in\mathcal M_v\\d(P)\ge r}}
\{1-\mathsf R_n(P,\phi)\}.
\]
For a nonempty bounded alternative set, this follows by inclusion; for an empty set, it follows from the zero convention and nonnegativity of the full-model error. Taking infima over full-null level-valid tests, and allowing all bounded-null level-valid tests on the left, proves
\[
B_n^{\mathrm b}(w,r)\le B_n(v,r).
\]
Together with the equality already established, this gives, for \(n\ge2\) and \(0<r<D_v\),
\[
B_n^{\mathrm{bin}}(w,r)=B_n^{\mathrm b}(w,r)\le B_n(v,r).
\]
% lean: CausalSmith.Stat.FinitepHomogeneityDensegamma.boundedTestingRisk_le_testingRisk

\item We establish strict positivity of the binary capped radius by a small cosine alternative. For \(t_{\mathrm{cos}}\in[0,1]\), define the deterministic law and its conversion by
\[
\begin{gathered}
P^{\mathrm{det}}_{\mathrm{cos},t_{\mathrm{cos}}}:
\quad X\sim\lambda,\qquad
\Pr(A=1\mid X=x)=\frac12,\qquad
Y=A\,\frac{t_{\mathrm{cos}}}{8}\cos(2\pi X),\\
P_{\mathrm{cos},t_{\mathrm{cos}}}
=\mathcal C_{\mathrm{bin}}
 \bigl(P^{\mathrm{det}}_{\mathrm{cos},t_{\mathrm{cos}}}\bigr).
\end{gathered}
\]
Their primitive functions are
\[
e_{P_{\mathrm{cos},t_{\mathrm{cos}}}}(x)=\frac12,\qquad
m_{0,P_{\mathrm{cos},t_{\mathrm{cos}}}}(x)=0,\qquad
\tau_{P_{\mathrm{cos},t_{\mathrm{cos}}}}(x)
=\frac{t_{\mathrm{cos}}}{8}\cos(2\pi x).
\]
The propensity and baseline are constant. For the effect, the cosine increment bound and \(|x-z|\le1\) imply
\[
\begin{split}
|\tau_{P_{\mathrm{cos},t_{\mathrm{cos}}}}(x)|
&\le\frac{t_{\mathrm{cos}}}{8}\le\frac12,\\
|\tau_{P_{\mathrm{cos},t_{\mathrm{cos}}}}(x)
 -\tau_{P_{\mathrm{cos},t_{\mathrm{cos}}}}(z)|
&\le t_{\mathrm{cos}}\frac{2\pi}{8}|x-z|
 \le t_{\mathrm{cos}}\,20|x-z|^\gamma
 \le20|x-z|^\gamma.
\end{split}
\]
The deterministic outcomes have absolute value at most \(1/8\), so their conditional second moments are at most one. These bounds verify membership of the deterministic law in
\(\mathcal M_{(2,\alpha,\beta,\gamma)}\); conversion then gives
\[
P_{\mathrm{cos},t_{\mathrm{cos}}}\in\mathcal M_w^{\mathrm{bin}},
\qquad
P_{\mathrm{cos},0}\in H_0^{\mathrm{bin}}(w).
\]
Moreover,
\[
\int\tau_{P_{\mathrm{cos},t_{\mathrm{cos}}}}\,d\lambda=0,\qquad
\int\tau_{P_{\mathrm{cos},t_{\mathrm{cos}}}}^2\,d\lambda
=\frac{t_{\mathrm{cos}}^2}{128},
\qquad
d(P_{\mathrm{cos},t_{\mathrm{cos}}})=t_{\mathrm{cos}}d_0.
\]
% lean: CausalSmith.Stat.FinitepHomogeneityDensegamma.smallCosineLaw_inBinaryModel
% lean: CausalSmith.Stat.FinitepHomogeneityDensegamma.smallCosineLaw_distance

For probability measures \(\mu_{\mathrm{tv}},\nu_{\mathrm{tv}}\) on a common measurable space \(\mathcal Z_{\mathrm{tv}}\), use
\[
\operatorname{TV}(\mu_{\mathrm{tv}},\nu_{\mathrm{tv}})
=
\sup_{E_{\mathrm{tv}}\text{ measurable in }\mathcal Z_{\mathrm{tv}}}
|\mu_{\mathrm{tv}}(E_{\mathrm{tv}})
 -\nu_{\mathrm{tv}}(E_{\mathrm{tv}})|.
\]
The conditional distribution of \((A,Y)\) under the cosine law is
\[
\begin{split}
\kappa_{\mathrm{cos},t_{\mathrm{cos}}}(x,\cdot)
={}&\frac{1+\tau_{P_{\mathrm{cos},t_{\mathrm{cos}}}}(x)}4
       \delta_{(1,1)}
+\frac{1-\tau_{P_{\mathrm{cos},t_{\mathrm{cos}}}}(x)}4
       \delta_{(1,-1)}\\
&+\frac14\delta_{(0,1)}+\frac14\delta_{(0,-1)}.
\end{split}
\]
For a measurable record event \(E_{\mathrm{tv}}\subseteq\mathcal O\), define its section by
\[
E_{\mathrm{tv},x}
=\{(a_{\mathrm{arm}},y):
  (x,a_{\mathrm{arm}},y)\in E_{\mathrm{tv}}\}.
\]
Subtracting the four displayed masses gives
\[
\begin{split}
&\kappa_{\mathrm{cos},0}(x,E_{\mathrm{tv},x})
-\kappa_{\mathrm{cos},t_{\mathrm{cos}}}(x,E_{\mathrm{tv},x})\\
&\quad=
-\frac{\tau_{P_{\mathrm{cos},t_{\mathrm{cos}}}}(x)}4
\left[
\mathbf1_{E_{\mathrm{tv},x}}(1,1)
-\mathbf1_{E_{\mathrm{tv},x}}(1,-1)
\right].
\end{split}
\]
The absolute difference of these two indicators is at most one. Therefore
\[
\left|
\kappa_{\mathrm{cos},0}(x,E_{\mathrm{tv},x})
-\kappa_{\mathrm{cos},t_{\mathrm{cos}}}(x,E_{\mathrm{tv},x})
\right|
\le
\frac{|\tau_{P_{\mathrm{cos},t_{\mathrm{cos}}}}(x)|}{4}
\le\frac{t_{\mathrm{cos}}}{16}.
\]
Integration over the common uniform design gives
\[
\begin{split}
|P_{\mathrm{cos},0}(E_{\mathrm{tv}})
-P_{\mathrm{cos},t_{\mathrm{cos}}}(E_{\mathrm{tv}})|
&\le
\int
\left|
\kappa_{\mathrm{cos},0}(x,E_{\mathrm{tv},x})
-\kappa_{\mathrm{cos},t_{\mathrm{cos}}}(x,E_{\mathrm{tv},x})
\right|\lambda(dx)\\
&\le\frac{t_{\mathrm{cos}}}{16},
\end{split}
\]
and hence
\[
\operatorname{TV}
(P_{\mathrm{cos},0},P_{\mathrm{cos},t_{\mathrm{cos}}})
\le\frac{t_{\mathrm{cos}}}{16}.
\]
% lean: CausalSmith.Stat.FinitepHomogeneityDensegamma.smallCosineLaw_tv

We use the expectation bound for any measurable
\(f_{\mathrm{tv}}:\mathcal Z_{\mathrm{tv}}\to[0,1]\). The layer-cake identity is
\[
\int f_{\mathrm{tv}}\,d\eta_{\mathrm{tv}}
=
\int_0^1
\eta_{\mathrm{tv}}
\{z_{\mathrm{tv}}:f_{\mathrm{tv}}(z_{\mathrm{tv}})\ge h_{\mathrm{tv}}\}
\,dh_{\mathrm{tv}},
\qquad
\eta_{\mathrm{tv}}\in\{\mu_{\mathrm{tv}},\nu_{\mathrm{tv}}\}.
\]
Subtracting these identities, taking absolute values, and bounding each event-probability difference gives
\[
\left|
\int f_{\mathrm{tv}}\,d\mu_{\mathrm{tv}}
-\int f_{\mathrm{tv}}\,d\nu_{\mathrm{tv}}
\right|
\le\operatorname{TV}(\mu_{\mathrm{tv}},\nu_{\mathrm{tv}}).
\]
Inserting the mixed product
\(P_{\mathrm{cos},0}\otimes
P_{\mathrm{cos},t_{\mathrm{cos}}}^{\otimes n_{\mathrm{prod}}}\)
between the two product laws, applying the triangle inequality, and integrating event sections gives
\[
\begin{split}
&\operatorname{TV}
\left(
P_{\mathrm{cos},0}^{\otimes(n_{\mathrm{prod}}+1)},
P_{\mathrm{cos},t_{\mathrm{cos}}}^{\otimes(n_{\mathrm{prod}}+1)}
\right)\\
&\quad\le
\operatorname{TV}(P_{\mathrm{cos},0},P_{\mathrm{cos},t_{\mathrm{cos}}})
+
\operatorname{TV}
\left(
P_{\mathrm{cos},0}^{\otimes n_{\mathrm{prod}}},
P_{\mathrm{cos},t_{\mathrm{cos}}}^{\otimes n_{\mathrm{prod}}}
\right),
\qquad n_{\mathrm{prod}}\in\mathbb N.
\end{split}
\]
For the first term, the integrated section probability is a \([0,1]\)-valued function, so the expectation bound applies; for the second, each section difference is bounded by the product total variation. At \(n_{\mathrm{prod}}=0\), both product laws are the point mass on the empty sample. Induction therefore yields
\[
\operatorname{TV}
\left(
P_{\mathrm{cos},0}^{\otimes n},
P_{\mathrm{cos},t_{\mathrm{cos}}}^{\otimes n}
\right)
\le
n\,\operatorname{TV}
(P_{\mathrm{cos},0},P_{\mathrm{cos},t_{\mathrm{cos}}})
\le\frac{nt_{\mathrm{cos}}}{16}.
\]
% lean: Causalean.Stat.Minimax.MomentMatchedMixture.tvDist_pi_iid_le

For the fixed \(n\ge2\), choose
\[
t_{\mathrm{small}}=\frac8{100n},
\qquad
0<t_{\mathrm{small}}\le1.
\]
The preceding product bound becomes
\[
\operatorname{TV}
\left(
P_{\mathrm{cos},0}^{\otimes n},
P_{\mathrm{cos},t_{\mathrm{small}}}^{\otimes n}
\right)
\le\frac{nt_{\mathrm{small}}}{16}
=\frac1{200}.
\]
For a binary-null level-valid test \(\phi\), define its seed average by
\[
\overline\phi(\boldsymbol o)
=\int_{[0,1]}\phi(\boldsymbol o,u_{\mathrm{seed}})
\,\lambda(du_{\mathrm{seed}}).
\]
It is measurable, satisfies \(0\le\overline\phi\le1\), and obeys
\[
\mathsf R_n(P,\phi)
=\int_{\mathcal O^n}\overline\phi\,dP^{\otimes n}.
\]
Applying the expectation bound to this average and using null validity gives
\[
\begin{split}
1-\mathsf R_n(P_{\mathrm{cos},t_{\mathrm{small}}},\phi)
&\ge
\frac9{10}
-\operatorname{TV}
\left(
P_{\mathrm{cos},0}^{\otimes n},
P_{\mathrm{cos},t_{\mathrm{small}}}^{\otimes n}
\right)\\
&\ge\frac{179}{200}>\frac1{10}.
\end{split}
\]
Whenever \(0<r\le t_{\mathrm{small}}d_0\), this cosine law belongs to the separated alternative set. Taking its error as a lower bound for the worst-case error, and then taking the infimum over level-valid tests, proves
\[
B_n^{\mathrm{bin}}(w,r)\ge\frac{179}{200}>\frac1{10}
\qquad(0<r\le t_{\mathrm{small}}d_0).
\]
% lean: CausalSmith.Stat.FinitepHomogeneityDensegamma.testingRiskOn_ge_of_two_point

Thus every successful separation in the binary capped-radius set exceeds \(t_{\mathrm{small}}d_0\). Its sentinel also satisfies
\[
t_{\mathrm{small}}d_0\le d_0\le D_w^{\mathrm{bin}}.
\]
The set is nonempty because it contains the sentinel. Taking its infimum, as defined in \cref{obj:def:binary-radius}, gives
\[
0<t_{\mathrm{small}}d_0
\le r_n^{*,\mathrm{bin}}(w).
\]
% lean: CausalSmith.Stat.FinitepHomogeneityDensegamma.binaryCriticalRadius_pos

\item Finally, the common cap and the risk equality give identical sets in the definitions of the binary and bounded radii:
\[
\begin{split}
&\left\{r\in(0,D_w^{\mathrm{bin}}):
 B_n^{\mathrm{bin}}(w,r)\le\frac1{10}\right\}
 \cup\{D_w^{\mathrm{bin}}\}\\
&\qquad=
\left\{r\in(0,D_w^{\mathrm b}):
 B_n^{\mathrm b}(w,r)\le\frac1{10}\right\}
 \cup\{D_w^{\mathrm b}\}.
\end{split}
\]
Their infima are therefore equal:
\[
r_n^{*,\mathrm{bin}}(w)=r_n^{*,\mathrm b}(w).
\]
The risk ordering also gives
\[
\begin{split}
&\left\{r\in(0,D_v):B_n(v,r)\le\frac1{10}\right\}\cup\{D_v\}\\
&\qquad\subseteq
\left\{r\in(0,D_v):B_n^{\mathrm b}(w,r)\le\frac1{10}\right\}
\cup\{D_v\}.
\end{split}
\]
Both sets contain \(D_v\) and are bounded below by zero, since \(D_v\ge d_0>0\). Taking infima reverses this inclusion. By
\cref{obj:def:bounded-radius,obj:def:critical-radius}, it follows that
\[
r_n^{*,\mathrm b}(w)\le r_n^*(v).
\]
Combining this with the positivity just proved establishes
\[
0<r_n^{*,\mathrm{bin}}(w)
=r_n^{*,\mathrm b}(w)\le r_n^*(v).
\]
The conversion properties and the identity for every \(n\in\mathbb N\) were established in the first step.
% lean: CausalSmith.Stat.FinitepHomogeneityDensegamma.cappedRadius_eq_of_risk_eq
% lean: CausalSmith.Stat.FinitepHomogeneityDensegamma.boundedCriticalRadius_le_criticalRadius
\end{enumerate}
\end{proof}
